\documentclass[11pt,letterpaper,oneside]{report}
\usepackage[T1]{fontenc}
\usepackage[utf8]{inputenc}
\usepackage{lmodern}
\usepackage[letterpaper,margin=0.92in,headheight=15pt]{geometry}
\usepackage{amsmath,amssymb,amsthm,mathtools}
\usepackage{microtype}
\usepackage[dvipsnames]{xcolor}
\usepackage{graphicx}
\usepackage{booktabs,longtable,array,tabularx}
\usepackage{enumitem}
\usepackage{fancyhdr}
\usepackage{titlesec}
\usepackage[most]{tcolorbox}
\usepackage{listings}
\usepackage{tikz}
\usetikzlibrary{positioning}
\usepackage{hyperref}
\hypersetup{colorlinks=true,linkcolor=MidnightBlue,urlcolor=RoyalBlue,pdfauthor={Independent Mathematics and Engineering Study Suite},pdftitle={Analytical Solution Methods Handbook}}
\definecolor{ink}{HTML}{162639}
\definecolor{navy}{HTML}{173B61}
\definecolor{blue}{HTML}{2874A6}
\definecolor{teal}{HTML}{007C83}
\definecolor{amber}{HTML}{8B5E16}
\definecolor{pale}{HTML}{EDF4F8}
\definecolor{paleteal}{HTML}{ECF8F5}
\definecolor{paleamber}{HTML}{FFF6E8}
\titleformat{\chapter}[display]{\normalfont\huge\bfseries\color{navy}}{\filleft\fontsize{43}{46}\selectfont\thechapter}{0.4ex}{\titlerule\vspace{1ex}}[\vspace{1ex}\titlerule]
\titleformat{\section}{\normalfont\Large\bfseries\color{navy}}{\thesection}{0.8em}{}
\titleformat{\subsection}{\normalfont\large\bfseries\color{teal}}{\thesubsection}{0.75em}{}
\pagestyle{fancy}\fancyhf{}\fancyhead[L]{\small\color{navy}Differential Equations Engineering Suite}\fancyhead[R]{\small\color{navy}Volume 3 of 8}\fancyfoot[C]{\thepage}\renewcommand{\headrulewidth}{0.3pt}
\setlist[itemize]{leftmargin=*,itemsep=3pt,topsep=5pt}\setlist[enumerate]{leftmargin=*,itemsep=4pt,topsep=5pt}
\setlength{\parskip}{5pt}\setlength{\parindent}{0pt}\setlength{\emergencystretch}{2em}
\newtheorem{theorem}{Theorem}[chapter]
\newtheorem{proposition}[theorem]{Proposition}
\newtheorem{lemma}[theorem]{Lemma}
\theoremstyle{definition}\newtheorem{definition}[theorem]{Definition}
\newtheorem{example}[theorem]{Worked example}
\theoremstyle{remark}\newtheorem{remark}[theorem]{Remark}
\newtcolorbox{bookground}[1][]{enhanced,breakable,colback=pale,colframe=blue,boxrule=0.6pt,title={TEXTBOOK ALIGNMENT},fonttitle=\sffamily\bfseries\small,#1}
\newtcolorbox{modern}[1][]{enhanced,breakable,colback=paleteal,colframe=teal,boxrule=0.6pt,title={MODERN COMPUTATIONAL EXTENSION},fonttitle=\sffamily\bfseries\small,#1}
\newtcolorbox{caution}[1][]{enhanced,breakable,colback=paleamber,colframe=amber,boxrule=0.6pt,title={ASSUMPTION / VALIDATION},fonttitle=\sffamily\bfseries\small,#1}
\newcommand{\dd}{\,\mathrm{d}}
\newcommand{\R}{\mathbb{R}}
\newcommand{\lap}{\mathcal{L}}
\newcommand{\norm}[1]{\left\lVert#1\right\rVert}
\newcommand{\abs}[1]{\left\lvert#1\right\rvert}
\lstdefinestyle{engpy}{language=Python,basicstyle=\ttfamily\footnotesize,keywordstyle=\color{RoyalBlue}\bfseries,commentstyle=\color{gray},stringstyle=\color{teal},backgroundcolor=\color{pale},frame=single,rulecolor=\color{blue!30},breaklines=true,columns=fullflexible,keepspaces=true,showstringspaces=false,tabsize=4,upquote=true}
\lstset{style=engpy}
\begin{document}
\begin{titlepage}
\pagecolor{navy}\color{white}
\vspace*{0.75in}
{\sffamily\bfseries\large MATHEMATICS \& ENGINEERING DOCUMENTATION SUITE\\[12pt]}
{\color{cyan!40}\rule{\textwidth}{2pt}}\\[1.0in]
{\sffamily\bfseries\fontsize{29}{36}\selectfont Analytical Solution Methods Handbook\par}
\vspace{0.45in}
{\Large Volume 3 of 8}\par\vspace{0.55in}
{\color{cyan!30}\rule{0.50\textwidth}{1pt}}\par\vspace{0.3in}
{\large Independent technical companion aligned with the\par}
{\large 13-chapter organization of Nagle, Saff \& Snider (7th ed.)\par}
\vfill
{\large Original derivations, worked examples, verification and engineering practice\par}\vspace{0.3in}
{\small Prepared October 2026 \quad\textbullet\quad pdf\LaTeX{} compatible\par}
\end{titlepage}
\nopagecolor\color{ink}\pagenumbering{roman}
\chapter*{Reader's guide and source boundaries}\addcontentsline{toc}{chapter}{Reader's guide and source boundaries}
This independent study resource is organized around \emph{Fundamentals of Differential Equations and Boundary Value Problems}, Seventh Edition, by R.~Kent Nagle, Edward~B. Saff, and Arthur~David Snider (Pearson, 2018; ISBN 978-0-321-97710-6). The book's chapter and section numbers are given as \emph{navigation references}, not as claims of reproduced text or proofs.

\begin{bookground}
\textbf{Textbook-grounded scope:} topics, terminology, and chapter/section organization corresponding to the cited textbook. All explanations, derivations, exercise prompts, and diagrams in this volume are newly authored and may emphasize different examples. Check the textbook itself for its exact hypotheses, proofs, exercises, or nomenclature in a particular section.
\end{bookground}
\begin{modern}
\textbf{Distinct extension layer:} Python/NumPy/SciPy workflows, engineering QA and reproducibility practices, automated convergence studies, diagnostic plots, and implementation guidance. These are \emph{not} represented as content supplied by the textbook. API calls are examples and should be verified against the installed package version.
\end{modern}
\begin{caution}
\textbf{Mathematical qualification:} formulas require their stated regularity, domain and boundary/initial conditions. Numerical agreement does not replace a proof; computational outputs are not certified error bounds unless explicitly established.
\end{caution}
\tableofcontents\clearpage\pagenumbering{arabic}
\chapter{The analytical solution workflow}
\begin{bookground}Method progression follows Chapters 2, 4, 6, 7 and 8. This handbook organizes the techniques into a reproducible decision workflow; it does not claim the original textbook orders every technique this way.\end{bookground}
\section{Preflight before calculation}
\begin{enumerate}
\item Classify order, linearity, autonomy and constant/variable coefficients.
\item Determine intervals on which coefficients and proposed transformations are defined.
\item Record initial or boundary data, smoothness and physical units.
\item Test recognizable structures before deploying more expensive methods.
\item Derive a general family, select constants, and \emph{verify by substitution}.
\item For BVPs, check existence/compatibility and possible nonuniqueness.
\end{enumerate}
\section{Separation and integrating factors}
If $y'=a(t)b(y)$, isolate and integrate away from equilibrium roots of $b$. If $y'+p y=q$, multiply by $\mu=e^{\int p\dd t}$. As a benchmark, solve $y'+2y=4$ with $y(0)=3$: $(e^{2t}y)'=4e^{2t}$, hence $y=2+e^{-2t}$.
\section{Exactness and substitutions}
For $M\dd x+N\dd y=0$, test $M_y=N_x$; find a scalar potential $F$. Common transformations are $v=y^{1-n}$ for Bernoulli equations and $y=vx$ when a right side depends on $y/x$. Transformation must be inverted within its valid domain, and excluded branches checked independently.
\begin{example}[Bernoulli equation]
Solve $y'+y=y^2$ on a branch with $y\ne0$. Set $v=1/y$, giving $v'-v=-1$. Thus $v=1+Ce^t$ and $y=1/(1+Ce^t)$ wherever the denominator is nonzero. The equilibrium $y\equiv0$ is also a solution; it was excluded by the substitution.
\end{example}
\chapter{Characteristic roots and forced equations}
\begin{bookground}Sections 4.2--4.6 and 6.2--6.4. The term \emph{auxiliary equation} is used in the textbook for constant-coefficient problems.\end{bookground}
\section{The exponential ansatz}
For $a_n D^ny+\cdots+a_0y=0$, substitute $e^{rt}$ to obtain $P(r)=0$. Distinct real roots contribute exponentials; conjugate pairs contribute real trigonometric modes; a root of multiplicity $m$ contributes $t^j e^{rt}$ for $0\le j<m$. Work on a domain where the operator is regular.
\section{Undetermined coefficients with resonance}
If $Ly=g(t)$ where $L$ has constant coefficients and $g$ is a finite sum of exponential-polynomial-sinusoidal terms, choose a trial family stable under differentiation. Multiply by $t^s$, where $s$ is the multiplicity of the associated characteristic root. Different forcing components can be treated by superposition.
\begin{example}[Polynomial forcing]
For $y''-y=2t$, try $y_p=At+B$. Then $-At-B=2t$, giving $A=-2$, $B=0$. Thus $y=C_1e^t+C_2e^{-t}-2t$.
\end{example}
\begin{example}[Repeated-root resonance]
For $(D-1)^2 y=e^t$, write $y=e^t v$. The shifted operator satisfies $(D-1)^2(e^t v)=e^t v''$, so $v''=1$ and $y_p=\tfrac12t^2e^t$. A $t^2$ factor is essential because $e^t$ and $te^t$ already span the homogeneous kernel.
\end{example}
\section{Variation of parameters in scalar form}
For normalized $y''+py'+qy=g$ with independent solutions $y_1,y_2$, the Wronskian $W=y_1y_2'-y_1'y_2$ gives
\[y_p=-y_1\int\frac{y_2g}{W}\dd t+y_2\int\frac{y_1g}{W}\dd t.\]
Choice of antiderivative constants changes only the homogeneous contribution. For equations with leading coefficient $a(t)$, normalize first. This technique accommodates forcings outside the undetermined-coefficients catalog.
\section{Cauchy--Euler equation}
On $t>0$, $t^2y''+at y'+by=0$ accepts $y=t^r$ and characteristic equation $r(r-1)+ar+b=0$. Real repeated roots produce a logarithmic second mode $t^r\log t$; complex roots $r=\alpha\pm i\beta$ give $t^\alpha\cos(\beta\log t)$ and $t^\alpha\sin(\beta\log t)$.
\chapter{Laplace transform: unilateral time-domain analysis}
\begin{bookground}Correspondence: Chapter 7, Sections 7.1--7.10: definition, properties, inverse transforms, IVPs, discontinuities, periodic functions, convolution, impulses and systems.\end{bookground}
\section{Definition and initial data}
For a function of suitable exponential order and piecewise regularity on $[0,\infty)$, define
\[F(s)=\lap\{f(t)\}(s)=\int_0^\infty e^{-st}f(t)\dd t.\]
When transforms exist, $\lap\{f'\}=sF(s)-f(0^+)$ and $\lap\{f''\}=s^2F(s)-sf(0^+)-f'(0^+)$. This makes Laplace transforms particularly effective for constant-coefficient IVPs with specified data.
\begin{example}[Transform an IVP]
Solve $y'+2y=1$, $y(0)=0$. Transform: $(s+2)Y(s)=1/s$, so $Y(s)=1/[s(s+2)]=\tfrac12(1/s-1/(s+2))$. Therefore $y(t)=\tfrac12(1-e^{-2t})$ for $t\ge0$.
\end{example}
\section{Shift theorem and unit-step forcing}
Let $u(t-a)$ denote the Heaviside step. For $a\ge0$,
\[\lap\{u(t-a)g(t-a)\}(s)=e^{-as}G(s).\]
The distinction between $u(t-a)g(t-a)$ and $u(t-a)g(t)$ is essential. For the latter, write $g(t)=h(t-a)$ with $h(\tau)=g(\tau+a)$ before applying the theorem.
\begin{example}[Switch-on load]
For $y'+y=u(t-1)$ with $y(0)=0$, transform to $(s+1)Y=e^{-s}/s$, giving
\[y(t)=u(t-1)\bigl(1-e^{-(t-1)}\bigr).\]
The solution is zero until activation and then approaches the forced equilibrium.
\end{example}
\section{Convolution and Green-type responses}
With $(f*g)(t)=\int_0^t f(t-\tau)g(\tau)\dd\tau$, the convolution theorem gives $\lap\{f*g\}=F(s)G(s)$. For $y'+ay=g(t)$ with $y(0)=0$,
\[y(t)=\int_0^t e^{-a(t-\tau)}g(\tau)\dd\tau.\]
This is the causal impulse-response representation, and it anticipates Green's functions for boundary-value problems.
\section{Impulse distributions and limitations}
The idealized Dirac delta is not an ordinary function. Operationally, its transform satisfies $\lap\{\delta(t-a)\}=e^{-as}$ for $a\ge0$. An impulse may create a jump in a state or its derivative according to the governing equation. Interpret the result in a distributional or weak sense; never approximate an impulse by treating $\delta$ as a finite-amplitude point measurement.
\chapter{Power series and Frobenius constructions}
\begin{bookground}Correspondence: Chapter 8, including \S8.3 power series, \S8.5 Cauchy--Euler, \S8.6 Frobenius and \S8.8 special functions.\end{bookground}
\section{Generate and check a recurrence}
For $y''+xy=0$, put $y=\sum_{n=0}^\infty a_n x^n$. Comparing $x^n$ gives
\[(n+2)(n+1)a_{n+2}+a_{n-1}=0\quad(n\ge1),\qquad2a_2=0.\]
Thus $a_2=0$, $a_3=-a_0/6$, $a_4=-a_1/12$, etc. The two free constants $a_0,a_1$ match the equation order. For a numerical approximation, truncate and differentiate the polynomial to evaluate a residual.
\section{Indicial equation at a regular singular point}
Writing $y=x^r\sum_{n\ge0}a_nx^n$ ($x>0$), the coefficient of the lowest power determines admissible $r$. For $x^2y''+xy'-\nu^2y=0$, the indicial equation is $r^2-\nu^2=0$ and solutions are $x^{\nu},x^{-\nu}$ when $\nu\ne0$; at $\nu=0$ the repeated root yields $1,\log x$.
\section{Beyond elementary closed forms}
The fact that an analytic solution uses Bessel, Airy, Legendre or Hermite functions is not a failure of the method: these functions encode families of solutions defined by differential equations and boundary behavior. Report normalization conventions and the numerical library convention used.
\chapter{Verification: residuals, data and singularities}
\section{The residual method}
For $Ly=f$, compute $r(t)=L\tilde y(t)-f(t)$. Check $r\equiv0$ when claiming an exact solution, and verify every initial/boundary condition separately. For an approximate series or numerical solution, report the residual norm on a chosen grid or function space; note that a small residual does not always imply a small error.
\section{Domain and branch checks}
Solutions such as $y=(t-c)^{-1}$ or $\log t$ require domain restrictions. Frobenius expansions may be centered at singular points but be valid on one-sided domains. Exponential transforms and Laplace integrals have regions of convergence; omitting them can turn a formal manipulation into an invalid claim.
\chapter{Modern computational extension: symbolic transforms and response checks}
\begin{modern}The Python example uses modern SymPy APIs; it is not a reproduction of Chapter 7 software instructions.\end{modern}
\begin{lstlisting}
import sympy as sp
t, s = sp.symbols('t s', positive=True)
y = sp.Function('y')
ode = sp.Eq(sp.diff(y(t), t) + 2*y(t), 1)
ans = sp.dsolve(ode, ics={y(0): 0})
assert sp.simplify(ans.rhs.subs(t, 0)) == 0
res = sp.simplify(sp.diff(ans.rhs, t) + 2*ans.rhs - 1)
assert res == 0
F = sp.laplace_transform(ans.rhs, t, s, noconds=True)
expected = 1 / (s * (s + 2))
assert sp.simplify(F - expected) == 0
\end{lstlisting}
A robust engineering check should also compare symbolic and independent high-accuracy numerical solutions, then vary initial values and parameters. Validate causal shifts by examining values immediately before and after the activation time.
\chapter{Exercises and independently authored solutions}
\section{Exercises}
\begin{enumerate}
\item Solve $y'+4y=12$, $y(0)=-1$.
\item Find the homogeneous solution of $y''+6y'+9y=0$.
\item Solve $y''+y=0$, $y(0)=2$, $y'(0)=-1$ using a transform or characteristic roots.
\item Find the convolution integral solution of $y'+3y=g(t)$, $y(0)=0$.
\item Derive $y_p$ for $y''+9y=\cos 3t$.
\item Derive the indicial exponents for $x^2y''+3xy'+y=0$.
\end{enumerate}
\section{Answer sketches}
\begin{enumerate}
\item $y=3-4e^{-4t}$.
\item $y=(C_1+C_2 t)e^{-3t}$.
\item $y=2\cos t-\sin t$.
\item $y(t)=\int_0^t e^{-3(t-\tau)}g(\tau)\dd\tau$.
\item $y_p=\tfrac16t\sin3t$, since $(D^2+9)(t\sin3t)=6\cos3t$.
\item $r(r-1)+3r+1=(r+1)^2$, so $y=x^{-1}(C_1+C_2\log x)$ for $x>0$.
\end{enumerate}

\appendix
\chapter{References and reproducibility}
\begin{enumerate}
\item R.~K. Nagle, E.~B. Saff, and A.~D. Snider, \emph{Fundamentals of Differential Equations and Boundary Value Problems}, 7th ed., Pearson, 2018, ISBN 978-0-321-97710-6. Primary source for the \emph{topic map and terminology}; cited locally by chapter and section.
\item SciPy project, \emph{solve\_ivp} reference, \url{https://docs.scipy.org/doc/scipy/reference/generated/scipy.integrate.solve_ivp.html} (modern extension).
\item SciPy project, \emph{solve\_bvp} reference, \url{https://docs.scipy.org/doc/scipy/reference/generated/scipy.integrate.solve_bvp.html} (modern extension).
\item SciPy project, \emph{eigsh} reference, \url{https://docs.scipy.org/doc/scipy/reference/generated/scipy.sparse.linalg.eigsh.html} (modern extension).
\item Python, NumPy, SciPy, and SymPy project documentation; computational checks require recording package versions and tolerances.
\end{enumerate}
\section*{Verification record template}
\begin{tabularx}{\textwidth}{>{\bfseries}p{0.30\textwidth}X}\toprule
Problem identifier & \rule{0pt}{15pt}\ \\ \midrule
ODE/PDE and domain & \\[8pt]
Initial/boundary conditions & \\[8pt]
Units and parameter values & \\[8pt]
Analytic residual & \\ [8pt]
Grid/step refinement & \\ [8pt]
Independent solver comparison & \\ [8pt]
Version and reproducibility notes & \\ [8pt]\bottomrule
\end{tabularx}
\end{document}
